\documentclass[12pt]{article}
\usepackage[utf8]{inputenc}
\usepackage[T1]{fontenc}
\usepackage{hyperref}
\title{Lachix\'io Zapotec Language Resources for K-12 Curriculum Development}
\author{Compiled by Claude Sonnet 5.5 (claude-sonnet-5-5) for Andr\'as Kornai}
\date{}
\begin{document}
\maketitle

\section{Introduction}
Lachix\'io Zapotec (ISO 639-3 zpl, Glottocode lach1249) is an Otomanguean language of Oaxaca, Mexico, with about 3,000 speakers. Searches found very little material aimed at schools. What exists is a lexical database and community dictionary project, an archival wordlist recording, and religious audio. No textbooks, primers, or school programmes were found.

\section{Primary Educational Resources}
No primers, textbooks or school curricula for Lachix\'io Zapotec turned up in search results. This is a real gap that curriculum work would have to fill from scratch.

\section{Linguistic and Reference Materials}
\begin{itemize}
\item Zapotec and Chatino project site (University of Virginia). It describes work to turn the Lachix\'io Zapotec lexical database, built by Sicoli with speakers from Santa Mar\'ia and San Vicente Lachix\'io, into an online collaborative community dictionary using the open-source Mukurtu software. The search results do not show whether the dictionary is live. \url{https://zapotec-chatino.iath.virginia.edu/}
\item Glottolog entry for the language. \url{https://glottolog.org/resource/languoid/iso/zpl}
\item Glosbe pages for Lachix\'io Zapotec. These are crowd-sourced, and I could not confirm what entries they hold; treat them with caution. \url{https://glosbe.com/lib/zpl}
\end{itemize}

\section{Religious and Community Resources}
\begin{itemize}
\item Global Recordings Network page for Lachix\'io Zapotec, listing evangelistic and Bible-teaching audio, and linking to a New Testament recording from Faith Comes By Hearing. \url{https://globalrecordings.net/en/language/zpl}
\end{itemize}

\section{Digital and Audio Resources}
\begin{itemize}
\item AILLA (University of Texas at Austin) archive item ``Elicitation of words from the dictionary'', a 2007 recording session from Santa Mar\'ia Lachix\'io with audio files and a wordlist. \url{https://islandora-ailla.lib.utexas.edu/islandora/object/ailla%3A136477}
\end{itemize}

\section{Curriculum Development Context}
Children in the Lachix\'io towns are reportedly raised with Zapotec and learn Spanish early, and the language is not known to be taught in schools. The most useful starting points are the Sicoli lexical database and community dictionary project and the AILLA recordings. A curriculum would need new materials in the language, an agreed orthography, and work with community speakers.

\end{document}